
\vspace{-0.3cm}
\section{Effects of {\color{cc1}optimization} on attention Sink.}
\vspace{-0.25cm}
\label{sec4}


We pre-train a series of LLaMA models to conduct our experiments, based on the repos~\citep{zhang2024tinyllama,liu2024regmix}. Due to the intractability of replicating LLaMA pre-training, we design small-sized models. Following \citet{liu2024regmix}, we set hidden dimension $d=\textrm{768}$, block number $L=\textrm{10}$, head number $H=\textrm{8}$, intermediate size of FFN as 1536, resulting in approximately 60M parameters except for word embeddings and unembeddings. We keep the other design the same as LLaMA2 models, including Rotary~\citep{su2024roformer}, pre-norm structure, RMSNorm~\citep{zhang2019root} as LN, SwiGLU activation~\citep{shazeer2020glu} in FFN, etc.\looseness=-1

For data distribution, we sample 5B tokens from the Pile dataset~\citep{gao2020pile}. We set the context length to 2048 tokens, the batch size to 1M tokens, and the training step to 20k (including 100 steps for warm-up). We adopt a learning rate of 4e-4 with cosine scheduling. The optimizer is AdamW~\citep{loshchilov2017decoupled} with a weight decay ratio of 0.1. We use the Pile-CC validation loss~\citep{gao2020pile,liu2024regmix} to measure the model performance and sample 100 sequences with $T=64$ (no BOS token) out of training data to measure the metric $\textrm{Sink}_k^\epsilon$ with $\epsilon=\textrm{0.3}$.  
% Validation is conducted on Pile-CC set, which contains about 200M tokens. 

% \subsection{Pre-Training hyper-parameters}



\begin{table}[t]
\vspace{-1.1cm}
\caption{Larger weight decay ratios tend to induce more attention sink heads in LMs. But much larger values hurt the model performance and attention sink disappears.}
\vspace{-.4cm}
\begin{center}
\begin{tabular}{l|cccccccc}
\toprule
$\gamma$ & 0.0 & 0.001 & 0.01 & 0.1 & 0.5 & 1.0 & 2.0 & 5.0\\
\midrule
$\textrm{Sink}_1^{\epsilon}$(\%)  & 15.20 & 15.39 & 15.23 & 18.18 & 41.08 & 37.71 & 6.13 &  0.01\\
valid loss & \,\,\,3.72 & \,\,\,3.72 & \,\,\,3.72 & \,\,\,3.73 & \,\,\,3.80  & \,\,\,3.90 &  4.23 & 5.24\\
\bottomrule
\end{tabular}
\end{center}
\vspace{-.5cm}
\label{decay}
\end{table}


% \begin{table}[t]
% \vspace{-1.1cm}
% \caption{(\emph{Left}) With a smaller learning rate, attention sink tends to be less obvious. However, a much smaller learning rate leads to suboptimal optimization, thus no attention sink. (\emph{Right}) Batch size has no effect on the emergence of attention sink.\looseness=-1}
% \vspace{-.2cm}
% \begin{minipage}[t]{0.49\linewidth}
% \begin{center}
% \begin{tabular}{l|cccc}
% \toprule
% lr & 1e-3 & 4e-4 & 1e-4 & 1e-5 \\
% \midrule
% $\textrm{Sink}_1^{\epsilon}$(\%) & 33.04 & 18.18 & 4.13 & 0.19\\
% valid loss & \,\,\,3.69 & \,\,\,3.73 & 3.78 & 4.69  \\
% \bottomrule
% \end{tabular}
% \end{center}
% \end{minipage}
% \hfill
% \begin{minipage}[t]{0.49\linewidth}
% \begin{center}
% \begin{tabular}{l|cccc}
% \toprule
% batch size & 0.25M & 0.5M & 1M & 2M \\
% \midrule
% $\textrm{Sink}_1^{\epsilon}$(\%) & 16.45 & 17.19 & 18.18 & 16.19 \\
% valid loss & \,\,\,3.92 & \,\,\,3.78 & \,\,\,3.73 & \,\,\,3.68 \\
% \bottomrule
% \end{tabular}
% \end{center}
% \end{minipage}
% \vspace{-.2cm}
% \label{lr_batch}
% \end{table}



% \begin{table}[htbp]
% \begin{center}
% \begin{tabular}{l|cccc}
% \toprule
% learning rate & 1e-3 & 4e-4 & 1e-4 & 1e-5 \\
% \midrule
% $\textrm{Sink}_1^{\epsilon}$(\%) & 33.04 & 18.18 & 4.13 & 0.19\\
% Valid loss & 3.69 & 3.73 & 3.78 & 5.36  \\
% \bottomrule
% \end{tabular}
% \end{center}
% \vspace{-.2cm}
% \caption{Learning rate.}
% \vspace{-.2cm}
% \label{lr}
% \end{table}


% \begin{table}[t]
% \begin{center}
% \begin{tabular}{l|cccc}
% \toprule
% Batch size & 0.25M & 0.5M & 1M & 2M \\
% \midrule
% $\textrm{Sink}_1^{\epsilon}$(\%) & 16.45 & 17.19 & 18.18 & 16.19 \\
% Valid loss & 3.92 & 3.78 & 3.73 & 3.68 \\
% \bottomrule
% \end{tabular}
% \end{center}
% \vspace{-.2cm}
% \caption{Batch size.}
% \vspace{-.2cm}
% \label{bs}
% \end{table}

%\subsection{Pre-training process and training hyper-parameters}

\textbf{Optimization steps.} As visualized in Figure~\ref{model}(\emph{Middle}), under our default setup, attention sink emerges after certain optimization steps, e.g., between 1k and 2k steps. With the progression of pre-training, attention sink becomes more obvious. 

\textbf{Learning rate.} With a smaller learning rate, it takes longer training steps to lower training loss, as present in Figure~\ref{model}(\emph{Right}). Meanwhile, the emergence of attention sink is also delayed. Besides, as shown in Table~\ref{lr}, we also find that a smaller learning rate results in LMs with less obvious attention sink, even if we compensate for more training steps.  But further decreasing learning rate significantly affects the optimization and model performance, thus affecting the emergence of attention sink.\looseness=-1

\textbf{Batch size.} In Table~\ref{batch_fix}(\emph{Left}), we find that only modifying batch size has no effects on attention sink.\looseness=-1

\vspace{-0.1cm}
\begin{abox} % are we tired of takeaway boxes? :>
    \looseness -1 \textbf{Takeaways:} 1.~Attention sink emerges after LMs are trained effectively. 2.~Attention sink appears less obvious in LMs trained with small learning rates.% at test time.
\end{abox}

% \textbf{Optimizer}


% \section{The Emergence of Attention Sink during the Inference}

% \begin{figure*}[t]
% \centering
% \subfigure{\includegraphics[width=0.23\textwidth]{figures/visuals/Llama-2-7b-hf/token_len200_probe3_first10.png}}
% \subfigure{\includegraphics[width=0.23\textwidth]{figures/visuals/Llama-2-7b-chat-hf/token_len200_probe3_first10.png}}
% \subfigure{\includegraphics[width=0.23\textwidth]{figures/visuals/Llama-2-13b-hf/token_len200_probe3_first10.png}}
% \subfigure{\includegraphics[width=0.23\textwidth]{figures/visuals/Llama-2-13b-chat-hf/token_len200_probe3_first10.png}}
% \subfigure{\includegraphics[width=0.23\textwidth]{figures/visuals/Meta-Llama-3-8B/token_len200_probe3_first10.png}}
% \subfigure{\includegraphics[width=0.23\textwidth]{figures/visuals/Meta-Llama-3-8B-Instruct/token_len200_probe3_first10.png}}
% \subfigure{\includegraphics[width=0.23\textwidth]{figures/visuals/Meta-Llama-3.1-8B/token_len200_probe3_first10.png}}
% \subfigure{\includegraphics[width=0.23\textwidth]{figures/visuals/Meta-Llama-3.1-8B-Instruct/token_len200_probe3_first10.png}}
% \subfigure{\includegraphics[width=0.23\textwidth]{figures/visuals/gpt2-small/token_len200_probe3_first10.png}}
% \subfigure{\includegraphics[width=0.23\textwidth]{figures/visuals/gpt2-medium/token_len200_probe3_first10.png}}
% \subfigure{\includegraphics[width=0.23\textwidth]{figures/visuals/gpt2-large/token_len200_probe3_first10.png}}
% \subfigure{\includegraphics[width=0.23\textwidth]{figures/visuals/gpt2-xl/token_len200_probe3_first10.png}}
% \subfigure{\includegraphics[width=0.23\textwidth]{figures/visuals/opt-350m/token_len200_probe3_first10.png}}
% \subfigure{\includegraphics[width=0.23\textwidth]{figures/visuals/opt-1.3b/token_len200_probe3_first10.png}}
% \subfigure{\includegraphics[width=0.23\textwidth]{figures/visuals/opt-2.7b/token_len200_probe3_first10.png}}
% \subfigure{\includegraphics[width=0.23\textwidth]{figures/visuals/opt-6.7b/token_len200_probe3_first10.png}}
% \subfigure{\includegraphics[width=0.23\textwidth]{figures/visuals/pythia-410m/token_len200_probe3_first10.png}}
% \subfigure{\includegraphics[width=0.23\textwidth]{figures/visuals/pythia-1b/token_len200_probe3_first10.png}}
% \subfigure{\includegraphics[width=0.23\textwidth]{figures/visuals/pythia-2.8b/token_len200_probe3_first10.png}}
% \subfigure{\includegraphics[width=0.23\textwidth]{figures/visuals/pythia-6.9b/token_len200_probe3_first10.png}}
% \caption{Attention distribution of token positions across various models.
% }
% \vspace{-0.35cm}
% \label{fig:attention_dist}
% \end{figure*}


% In the last section, we have shown that the initial token attention sink is a universal scenario. Next, we would like to explore in which circumstances this scenario emerges during the inference.  


% \subsection{Repeating Input Tokens}


% \begin{corollary}\label{corollary2}
% Suppose the first $T'$ input tokens are the same $\boldsymbol{x}_t=\boldsymbol{x}, 1 \leq t \leq T', T'<T$. If $f_{\theta}$ has the relative positional embeddings, ALiBi or Rotary as the positional embeddings,  all hidden states for token position $1 \leq t \leq T'$ are the same. The relative distance defines the distributions of attention weights 
% for these token positions. For learnable positional embeddings and absolute positional embeddings, the above conclusions do not hold.
% \end{corollary}


% \subsection{Modifying Position IDs}

% \textbf{Replacing Position IDs.} Replace the first position with another position.

% \textbf{Swapping Position IDs.} Swap two positions. 


% \subsection{Modifying Attention Mask}

% In this section, we first identify an ideal case where no attention sink is theoretically guaranteed. 



% \textbf{Conclusion:} For a decoder-only Transformer without position embedding, if all the input tokens are the same, the hidden states for each token position $t$ are the 
% same. The attention weights for each token position $t$ are uniform distributions without attention sinks.

